% !TeX root = ../main.tex

\chapter{Conclusion and Outlook}\label{chapter:conclusion-outlook}

\section{Conclusion}

This thesis investigates the building energy management problem of multiple prosumers in a low-voltage distribution network. Each prosumer is equipped with household loads, PV generation, stationary battery energy storage system, and residential EV charging demand. The objective is to coordinate these resources under dynamic electricity prices while considering EV mobility requirements and distribution network safety.
% 本文研究低压配电网中多个产消者的协调能源管理问题。每个产消者均配置家庭用电负荷、光伏发电、固定式电池储能系统以及住宅电动汽车充电需求。研究目标是在动态电价条件下协调这些资源，同时兼顾电动汽车出行需求和配电网安全。

A complete modeling, control, and evaluation framework is established. The system model includes household loads, PV generation, stationary batteries, EV charging, electricity prices, net load calculation, and power flow constraints. The battery and EV models describe the main physical constraints, including SoC dynamics, power limits, EV connection windows, arrival SoC, and departure requirements. Based on this model, the proposed control framework combines LSTM forecasting, MADRL actor--critic control, EV constraint handling, and optional grid safety projection.
% 本文建立了一套完整的建模、控制与评估框架。系统模型包括家庭负荷和热泵负荷、光伏发电、固定式电池、电动汽车充电、电价、净负荷计算以及潮流约束。电池与电动汽车模型描述了主要物理约束，包括 SoC 动态、功率限制、电动汽车连接窗口、到达时 SoC 和离开要求。在此模型基础上，所提出的控制框架结合了 LSTM 预测、MADRL 演员—评论家控制、电动汽车约束处理以及可选的电网安全投影。

The final results show that the LSTM model can capture the main trend of electricity price variations and provides useful forecast information for the controller, and the MADRL training results indicate that a sufficiently long training process is necessary for learning stable scheduling behavior in the presence of battery SoC, EV SoC, electricity price, and grid feedback.
% 结果表明，LSTM 模型能够捕捉电价变化的主要趋势，并为控制器提供有用的预测信息。MADRL 训练结果显示，在同时考虑电池 SoC、电动汽车 SoC、电价和电网反馈的情况下，需要足够长的训练过程才能学习到稳定的调度行为。

This thesis gives different kinds of controlling methods for EV departure requirements. For EV soft constraints, reward-only control can learn battery arbitrage behavior but can not well satisfy the EV departure SoC requirement. The Progress penalty improves EV feasibility and provides a better balance between EV charging requirements and economic performance than simple cost-weight adjustment. For EV hard constraints, this method can ensure the EV departure requirement at the action execution layer and improve mobility feasibility. However, EV hard handling alone is not sufficient to guarantee grid safety.
% 对于电动汽车软约束，仅依赖奖励函数的控制可以学习电池套利行为，但无法可靠满足电动汽车离开时的 SoC 要求。Progress 惩罚提高了电动汽车可行性，并且与简单调整成本权重相比，能够在电动汽车充电需求和经济性能之间取得更好的平衡。对于电动汽车硬约束，在动作执行层强制满足车辆离开要求可以提高出行可行性。然而，仅采用电动汽车硬约束处理仍不足以保证电网安全。

Grid safety projection improves distribution network operation by reducing safety violations. At the same time, it may reduce battery arbitrage performance because the executed actions can differ from the raw actions generated by the MADRL actor. Therefore, the results reveal a clear trade-off between economic performance and grid safety when action projection is applied.
% 电网安全投影通过减少安全约束违反改善配电网运行。同时，由于最终执行动作可能不同于 MADRL 演员生成的原始动作，该机制也可能降低电池套利性能。因此，结果表明在采用动作投影时，经济性能与电网安全之间存在明显的权衡关系。

% Modified: use the readable MADRL variant names defined in Table~\ref{tab:madrl-name-mapping}.
Compared with the rule-based battery controller, MADRL learns more flexible price-aware and resource-coupling policies. Compared with MPC, MADRL still shows weaker economic performance, but it provides fast decentralized execution after training. Overall, the results demonstrate that MADRL is a promising approach for multi-prosumer energy management in low-voltage distribution networks. In the studied scenario, Base Hard is more suitable when economic performance is prioritized, while Projected Hard is more robust when EV feasibility and grid safety are emphasized.
% 与基于规则的电池控制器相比，MADRL 能够学习更灵活的电价感知策略和资源耦合策略。与 MPC 相比，MADRL 的经济性能仍然较弱，但训练完成后可以实现快速的分散执行。总体而言，结果证明 MADRL 是一种适用于低压配电网多产消者能源管理的有前景的方法。在本文研究场景中，当优先考虑经济性能时，\texttt{Base Hard} 更为合适；当重点关注电动汽车可行性和电网安全时，\texttt{Projected Hard} 更具鲁棒性。

\section{Outlook}

Although this thesis has established a relatively complete framework for multi-prosumer modeling, control, and evaluation, and has systematically compared the performance of MADRL, MPC, and the rule-based battery controller, some directions remain worthy of further investigation.
% 尽管本文已经建立了完整的多产消者建模、控制与评估框架，并系统比较了 MADRL、MPC 和基于规则的电池控制器的性能，但仍有若干方向值得进一步研究。

\textbf{Projection-aware MADRL training:} The experiments in this thesis demonstrate that grid safety projection can significantly improve distribution network safety but may also weaken battery arbitrage. The fundamental reason is that the raw actions generated by the actor are inconsistent with the projected actions ultimately executed by the environment. Future work could incorporate the projection mechanism more directly into the training process. For example, the critic could learn the effects of the projected actions on the reward and the next state, or the actor could directly receive feedback based on the projected actions during training. This would enable the policy to adapt to the ultimately executed actions during learning, thereby reducing the negative impact of action projection on economic performance.
% \textbf{投影感知的 MADRL 训练：}本文实验表明，电网安全投影能够显著提高配电网安全性，但也可能削弱电池套利。其根本原因在于演员生成的原始动作与环境最终执行的投影后动作不一致。未来研究可以将投影机制更直接地纳入训练过程。例如，评论家可以学习投影后动作对奖励和下一状态的影响，或者演员在训练期间直接接收基于投影后动作的反馈。这样能够使策略在学习过程中适应最终执行动作，从而减小动作投影对经济性能的不利影响。

\textbf{Larger-scale and more heterogeneous distribution network scenarios:} This thesis considers a low-voltage distribution network consisting of \(N=3\) prosumers, which is suitable for analyzing the fundamental coupling among EVs, stationary batteries, PV systems, and distribution network safety. However, practical distribution networks may include a larger number of households, more complex feeder topologies, PV systems and batteries with different capacities, diverse EV arrival and departure patterns, and a higher proportion of uncertain loads. Future work could extend the proposed framework to larger SimBench networks or real-world distribution network datasets in order to evaluate the scalability and robustness of MADRL under more complex operating conditions.
% \textbf{更大规模且异质性更强的配电网场景：}本文考虑由 \(N=3\) 个产消者组成的低压配电网，适合分析电动汽车、固定式电池、光伏系统和配电网安全之间的基本耦合关系。然而，实际配电网可能包含更多家庭、更复杂的馈线拓扑、容量不同的光伏系统和电池、多样化的电动汽车到达与离开模式，以及更高比例的不确定负荷。未来可以将所提出的框架扩展到更大规模的 SimBench 网络或真实配电网数据集，以评估 MADRL 在更复杂运行条件下的可扩展性和鲁棒性。

\textbf{Real-time deployment and hybrid control methods:} MPC demonstrates superior economic performance, but its computational burden increases as the number of prosumers and the complexity of constraints grow. In contrast, MADRL offers fast online execution after training but requires extensive simulations and carefully designed reward functions during training. Future work could systematically compare the computation time of MADRL, Local MPC, and ADMM-MPC, and further explore hybrid control approaches. For example, MPC could be used to generate expert trajectories for imitation learning or to warm-start the MADRL policy, thereby combining the optimization capability of MPC with the fast execution of MADRL.
% \textbf{实时部署与混合控制方法：}MPC 表现出更优的经济性能，但其计算负担会随着产消者数量和约束复杂度的增加而上升。相比之下，MADRL 在完成训练后可以快速在线执行，但训练期间需要大量仿真以及经过谨慎设计的奖励函数。未来可以系统比较 MADRL、Local MPC 和 ADMM-MPC 的计算时间，并进一步探索混合控制方法。例如，可以利用 MPC 为模仿学习生成专家轨迹，或者对 MADRL 策略进行热启动，从而结合 MPC 的优化能力与 MADRL 的快速执行能力。

In summary, future work should focus on improving the awareness of MADRL with respect to execution-layer correction mechanisms, enhancing the coordination among EV feasibility, economic performance, and distribution network safety, and validating its generalization capability in larger-scale and more heterogeneous prosumer networks. These extensions will further improve the practical applicability of MADRL-based energy management for residential low-voltage distribution networks.
% 综上，未来研究应重点提高 MADRL 对执行层修正机制的感知能力，加强电动汽车可行性、经济性能和配电网安全之间的协调，并在规模更大、异质性更强的产消者网络中验证其泛化能力。这些扩展将进一步提高基于 MADRL 的能源管理方法在住宅低压配电网中的实际适用性。
